\documentclass[11pt,a4paper]{article}

\usepackage[T1]{fontenc}
\usepackage[utf8]{inputenc}
\usepackage[spanish,es-nodecimaldot]{babel}
\usepackage{lmodern}
\usepackage{microtype}
\usepackage{amsmath,amssymb,mathtools}
\usepackage{geometry}
\usepackage{enumitem}
\usepackage{booktabs}
\usepackage{tabularx}
\usepackage{xcolor}
\usepackage[hidelinks]{hyperref}

\geometry{margin=2.5cm}
\setlength{\parindent}{0pt}
\setlength{\parskip}{0.45em}
\setlist{itemsep=0.25em,topsep=0.3em}
\renewcommand{\arraystretch}{1.25}
\setcounter{secnumdepth}{0}

\newcommand{\argmax}{\operatorname*{arg\,max}}
\newcommand{\concepto}[1]{%
  \begin{center}
    \fcolorbox{black}{gray!8}{\parbox{0.88\linewidth}{\centering #1}}
  \end{center}%
}

\begin{document}

{\LARGE\bfseries Multi-Armed Bandits\par}
{\large Resumen del segundo capítulo de \textit{Reinforcement Learning: An Introduction}, de
Richard S. Sutton y Andrew G. Barto.\par}
\vspace{0.7em}
{\normalsize\bfseries INTRODUCCIÓN TEÓRICO-PRÁCTICA AL APRENDIZAJE POR REFUERZO\par}
\vspace{0.35em}
\textbf{Profesor:} Matias Bilkis
\hfill
\textbf{Cuatrimestre:} Segundo cuatrimestre 2026

\vspace{0.25em}
\textbf{Integrantes}

\begin{tabularx}{\textwidth}{@{}X@{\hspace{2em}}X@{}}
Agustín Calo \hfill 390/23 & Paz Lavenia \hfill 64/23 \\
Rocco Gasparini \hfill 102/24 & Silvano Picard \hfill 637/24 \\
Maria Delfina Kiss \hfill 72/23 & Ingrid von Foerster \hfill 150/23
\end{tabularx}

\vspace{0.35em}
\hrule

\section*{Idea general del capítulo}

La idea central del \textbf{aprendizaje por refuerzo}
(\textit{Reinforcement Learning}, RL) es bastante simple: el agente no recibe la
respuesta correcta de antemano. Prueba una acción y, a cambio, obtiene una
\textbf{evaluación numérica de esa acción: la recompensa o \textit{reward}}.

Esto lo diferencia del aprendizaje supervisado:

El \textbf{feedback evaluativo} indica qué tan buena fue la acción realizada.
Esto diferencia al aprendizaje por refuerzo del aprendizaje supervisado, que
trabaja con \textbf{feedback instructivo}: ahí sí se indica cuál era la acción
o respuesta correcta.

Como consecuencia, en RL el agente tiene que \textbf{probar acciones mediante
prueba y error} para descubrir cuáles son mejores.

El problema de los \textit{multi-armed bandits} simplifica esta idea eliminando estados y consecuencias futuras. Solamente hay que elegir acciones y
observar recompensas. Así podemos enfocarnos en uno de los problemas
centrales de RL: la exploración frente a la explotación.

\section{Exploración frente a explotación}

Un multi-armed bandit, dadas $n$ acciones posibles:

\[
  a_1,a_2,\dots,a_n.
\]

En cada instante $t$:

\begin{enumerate}
  \item elegimos una acción $A_t$;
  \item recibimos una recompensa $R_t$;
  \item repetimos el proceso.
\end{enumerate}

Cada acción tiene una distribución de recompensas propia, y el objetivo es
\textbf{maximizar la recompensa total esperada a lo largo del tiempo}.

(El nombre viene de las máquinas tragamonedas: una máquina de una sola palanca
es un \textit{one-armed bandit}; aquí imaginamos una máquina con $n$ palancas.)

Además, cada acción $a$ tiene un valor verdadero $q(a)$, que representa su
\textbf{recompensa esperada}:

\begin{equation}
    q(a)=\mathbb{E}\!\left[R_t\mid A_t=a\right]
\end{equation}
  
Si conociéramos todos los valores $q(a)$, resolver el problema sería trivial:
elegiríamos siempre la acción con mayor valor. El problema es que
\textbf{no conocemos esos valores}, por lo que tenemos que estimarlos.


Suponiendo que tenemos estimaciones de los valores, denotadas por $Q_t(a)$.
La acción con mayor valor estimado se denomina \textit{greedy action}.

\subsection{Explotación}

Consiste en elegir la acción que actualmente creemos que es mejor:

\begin{equation}
    A_t=\argmax_a Q_t(a)
\end{equation}

Estamos aprovechando el conocimiento que ya tenemos. Su ventaja es que maximiza la
recompensa esperada \textbf{inmediatamente}; su problema es que podríamos estar
equivocados acerca de cuál es la mejor acción.

\subsection{Exploración}

Consiste en elegir una acción que actualmente no parece ser la mejor para
obtener más información sobre ella. Nos permite descubrir acciones mejores,
pero con el costo de posiblemente producir menos recompensa en el corto plazo.

%\concepto{$\text{Exploración}\quad\text{frente a}\quad\text{Explotación}$}

No podemos explorar y explotar simultáneamente con una misma elección.
Sacrificar recompensa ahora puede permitir obtener obtener una recompensa mucho mayor en el
futuro. La decisión óptima depende, entre otras cosas, de:

\begin{enumerate}
  \item las estimaciones actuales
  \item la incertidumbre
  \item la cantidad de pasos restantes
  \item la estabilidad del problema
\end{enumerate}
  
Se buscan \textbf{métodos simples que consigan un buen equilibrio}. Para \textbf{estimar cuánto vale cada acción}, la idea es hacer dos cosas:

\begin{enumerate}
  \item estimar el valor de cada acción;
  \item utilizar esas estimaciones para decidir qué acción elegir.
\end{enumerate}

Se distingue entre $q(a)$, el valor \textbf{real} de la acción, y $Q_t(a)$,
su valor \textbf{estimado} en el tiempo $t$.

\subsection{Promediar de recompensas como estimación de valor}

La forma más directa de estimar el valor de una acción es promediar todas las
recompensas que obtuvimos al elegirla. A esto se lo llama Simple-average method. Si elegimos $a$ un total de
$N_t(a)$ veces:

\begin{equation}
    Q_t(a)=\frac{R_1+R_2+\cdots+R_{N_t(a)}}{N_t(a)}.
\end{equation}
  
Si la acción todavía nunca fue elegida, puede asignarse un valor inicial como $Q_1(a)=0$.

En un problema \textbf{estacionario}, si probamos la acción infinitas veces,
entonces, por la ley de los grandes números, tenemos que $Q_t(a)\longrightarrow q(a)$.

Es decir que la estimación del valor de la acción $a$ se aproxima cada vez más a su valor verdadero a medida que aumenta la experiencia. La estrategia más simple es elegir siempre la acción con mayor valor estimado. A esto se lo denomina
\textbf{Greedy Action Selection}:

\begin{equation}
  A_t=\argmax_a Q_t(a) 
\end{equation}
  
Esta es una política \textit{greedy}. El problema es que únicamente explota:
si una acción tuvo mala suerte en sus primeras recompensas, podría parecer peor
de lo que realmente es y no volver a ser elegida jamás.

Una solución muy sencilla consiste en introducir exploración aleatoria: \texorpdfstring{$\varepsilon$}{epsilon}-greedy

\begin{itemize}
  \item con probabilidad $1-\varepsilon$, elegimos la acción \textit{greedy};
  \item con probabilidad $\varepsilon$, elegimos una acción aleatoriamente.
\end{itemize}

Por ejemplo, con $\varepsilon=0.1$, aproximadamente el 90\,\% de las decisiones
son de explotación y el 10\,\% son de exploración.

Como existe una probabilidad positiva de elegir cualquier acción, a largo plazo
todas las acciones son probadas repetidamente. Por lo tanto tenemos que $N_t(a)\longrightarrow\infty$ y las estimaciones terminan acercándose a sus valores verdaderos.

El libro compara estos algoritmos mediante un problema experimental con $n=10$
acciones. Los valores verdaderos $q(a)$ se generan con una distribución normal $\mathcal{N}(0, \sigma^2)$ con $\sigma^2 = 1$. Las recompensas también contienen ruido aleatorio. Los resultados se promedian sobre \textbf{2000 problemas diferentes durante 1000 pasos}. Se comparan:

\begin{itemize}
  \item \textit{greedy}: $\varepsilon=0$;
  \item $\varepsilon$-greedy con $\varepsilon=0.01$;
  \item $\varepsilon$-greedy con $\varepsilon=0.1$.
\end{itemize}

El método \textit{greedy} funciona ligeramente mejor al principio, pero después
queda estancado porque frecuentemente encuentra una acción subóptima y deja de
explorar. En el experimento encontró la acción óptima solamente en
aproximadamente \textbf{un tercio de los problemas}.

En cambio, $\varepsilon$-greedy continúa explorando y eventualmente descubre
mejores acciones. Con $\varepsilon=0.1$ se descubre rápidamente la mejor acción,
pero se sigue explorando bastante incluso después. Con $\varepsilon=0.01$ se
aprende más lentamente, pero a largo plazo se explota más la acción óptima.
Incluso se podría hacer que $\varepsilon\longrightarrow 0$ con el tiempo.

Cuanto mayor sea el ruido de las recompensas, más difícil es identificar la
mejor acción y mayor valor tiene la exploración. También es fundamental en
problemas \textbf{no estacionarios}, donde $q_t(a)$ puede cambiar con el tiempo y significa el valor verdadero de la acción $a$ en el instante $t$. Una acción que antes era mala puede convertirse posteriormente en la mejor. Por eso, el equilibrio entre exploración y explotación es una característica fundamental de RL.

Ahora bien, supongamos que queremos calcular

\begin{equation}
    Q_n=\frac{R_1+\cdots+R_n}{n}
\end{equation}

La forma más directa de hacer esta cuenta sería guardar \textbf{todas las
recompensas anteriores}. El problema es que cada vez necesitaríamos más memoria
y más cómputo. Por suerte, podemos actualizar el promedio de a poco. Si ya
tenemos $Q_k$ y recibimos $R_k$:

\begin{equation}
    \boxed{Q_{k+1}=Q_k+\frac{1}{k}\left(R_k-Q_k\right)}
\end{equation}
  
Así solamente necesitamos almacenar el valor actual $Q_k$ y el contador $k$. Y de esta
manera, podemos actualizar estimaciones sin guardar todo

%\section{Una regla que nos aparece}

Esta fórmula es importante porque tiene una estructura que se vuelve a ver
una y otra vez en aprendizaje por refuerzo:

\[
  \boxed{
  \text{Nueva estimación}
  =\text{Estimación anterior}
  +\text{Tamaño de paso} \times
  \left[\text{Objetivo}-\text{Estimación anterior}\right]
  }.
\]

De forma abreviada:

\[
  \boxed{
  \text{Nueva estimación}
  =\text{Estimación anterior}+\alpha\times\text{Error}
  },
\]

donde $\text{Error}=\text{Objetivo}-\text{Estimación anterior}$

La nueva estimación se mueve parcialmente hacia un \textbf{objetivo}
(\textit{target}). Esta estructura es una de las ideas más importantes del
capítulo porque reaparece constantemente en algoritmos posteriores.

Hasta ahora veníamos suponiendo que $q(a)$ era constante. Pero en muchos entornos se tiene que $q_t(a)\neq q_{t+1}(a)$. En esos casos, promediar \textbf{toda la historia} puede jugar en contra: las recompensas de hace mucho tiempo quizás ya no representen lo que está pasando
ahora.

Para dar mayor importancia a las observaciones recientes se utiliza la siguiente formula: 

\begin{equation}
    \boxed{Q_{k+1}=Q_k+\alpha\left(R_k-Q_k\right)}
    \qquad \text{ con } 0<\alpha\leq 1
\end{equation}
  
Al expandir la fórmula se obtiene

\begin{equation}
    Q_{k+1}=(1-\alpha)^kQ_1
  +\sum_{i=1}^{k}\alpha(1-\alpha)^{k-i}R_i.
\end{equation}

Las recompensas más recientes reciben mayor peso. El peso de una recompensa
antigua disminuye aproximadamente de forma exponencial: $\alpha(1-\alpha)^{k-i}$

Por eso se denomina \textbf{promedio ponderado exponencialmente por recencia}
(\textit{exponential recency-weighted average}).

Si $\alpha$ es grande, por ejemplo $\alpha=0.8$, la estimación cambia
rápidamente ante nueva información: aprende y olvida rápido, pero puede resultar
más ruidosa.

Si $\alpha$ es pequeño, por ejemplo $\alpha=0.01$, la estimación cambia
lentamente, es más estable y conserva información antigua durante más tiempo.

Si utilizamos una secuencia de tamaños de paso $\alpha_k(a)$, para garantizar
convergencia suelen exigirse las siguientes condiciones:

\begin{itemize}
    \item $\sum_{k=1}^{\infty}\alpha_k(a)=\infty$
    \item $\sum_{k=1}^{\infty}\alpha_k^2(a)<\infty$
\end{itemize}

%\[
%  \sum_{k=1}^{\infty}\alpha_k(a)=\infty \text{  y  }
%  \sum_{k=1}^{\infty}\alpha_k^2(a)<\infty.
%\]

La primera evita que el algoritmo deje de aprender demasiado pronto. La segunda
hace que los cambios eventualmente sean suficientemente pequeños como para
permitir la convergencia.

Por ejemplo, tomando $\alpha_k=\frac{1}{k}$ se tiene que cumple ambas condiciones. En cambio, un valor constante como $\alpha_k=\alpha$ no cumple la segunda. Esto \textbf{no es necesariamente malo}: en un problema
no estacionario queremos que $Q_t(a)$ siga cambiando para adaptarse a las
modificaciones del entorno. Por eso los tamaños de paso constantes son muy
útiles en RL.

Toda estimación necesita un punto de partida $Q_1(a)$. En vez de usar, por
ejemplo, $Q_1(a)=0$, podemos arrancar deliberadamente con valores demasiado
altos, como $Q_1(a)=5$ aunque las recompensas reales estén aproximadamente alrededor de cero. Estos se denominan \textbf{valores iniciales optimistas}.

Suponiendo que $Q_1(a)=5$ para todas las acciones. Elegimos una acción y
obtenemos $R=1$; entonces, su estimación disminuye. Las demás todavía tienen
$Q(a)=5$, por lo que parecen mejores y el agente pasa a probarlas. El proceso se
repite y el agente termina probando muchas acciones aun utilizando una política
puramente \textit{greedy}.

La ventaja es que se trata de una forma muy facil de generar
exploración y no requiere $\varepsilon>0$.

El problema es que la exploración producida es \textbf{temporal}. Una vez que
desaparece el optimismo inicial, el algoritmo deja de explorar. Por eso no
funciona especialmente bien en problemas no estacionarios: si el entorno cambia
después de mucho tiempo, el algoritmo no recibe un nuevo impulso para explorar.

El problema de $\varepsilon$-greedy es que, cuando explora, lo hace casi
\textbf{a ciegas}: elige al azar entre las acciones. Eso significa que puede
volver a probar acciones claramente malas y dejar otras, más prometedoras pero
poco conocidas, en segundo plano. UCB (\textit{Upper Confidence Bound}) intenta
explorar con un poco más de criterio.

\section{El Upper Confidence Bound (UCB)}

\subsection{Formulación}

Se elige

\begin{equation}
      \boxed{
  A_t=\argmax_a\left[
  Q_t(a)+c\sqrt{\frac{\ln t}{N_t(a)}}
  \right]
  }

\end{equation}

La expresión tiene dos partes:

\begin{enumerate}
  \item \textbf{Explotación:} $Q_t(a)$ favorece las acciones con mayor recompensa
        estimada.
  \item \textbf{Exploración:} $c\sqrt{\frac{\ln t}{N_t(a)}}$ favorece las
        acciones sobre las que tenemos mayor incertidumbre.
\end{enumerate}

$N_t(a)$ indica cuántas veces elegimos la acción. Si fue elegida pocas veces, $N_t(a)$ es pequeño y $\sqrt{\frac{\ln t}{N_t(a)}}$ es grande. Por lo tanto, UCB la considera atractiva para explorar. Cuando la
probamos muchas veces, $N_t(a)$ aumenta y el incentivo de exploración disminuye.

El parámetro $c>0$ controla cuánto se valora la exploración, por tanto:

\begin{itemize}
  \item un $c$ grande produce más exploración;
  \item un $c$ pequeño produce más explotación.
\end{itemize}

\subsection{Su idea detrás}

El UCB puede interpretarse como la elección de la acción con el mayor
\textbf{valor plausible optimista}:

\[
  \underbrace{Q_t(a)}_{\text{estimación}}
  +
  \underbrace{\text{incertidumbre}}_{\text{exploración}}.
\]

Una acción puede ser elegida porque creemos que es buena o porque todavía
existe una posibilidad razonable de que sea mucho mejor de lo que pensamos.

En el experimento del libro, UCB supera generalmente a $\varepsilon$-greedy.
Sin embargo, se señala que resulta más difícil aplicar directamente
estas ideas a problemas generales de RL, especialmente cuando existen:

\begin{itemize}
  \item problemas no estacionarios;
  \item espacios de estados grandes;
  \item aproximación de funciones.
\end{itemize}

Hasta ahora, el agente siempre enfrentaba mas o menos \textbf{la misma
situación}. Por eso, el problema era encontrar la mejor acción. En el
aprendizaje por refuerzo completo existen diferentes situaciones, de modo que
tenemos que aprender una relación situación $\longrightarrow$ acción, es decir, una \textbf{política}.

Supongamos que existen varios bandits diferentes y que cada vez que jugamos
aparece uno distinto. Si no sabemos cuál apareció, el problema parece simplemente
no estacionario. Pero imaginemos que recibimos una pista o \textbf{contexto}; por
ejemplo:

\begin{itemize}
  \item pantalla roja $\longrightarrow$ bandit A;
  \item pantalla verde $\longrightarrow$ bandit B.
\end{itemize}

Entonces podemos aprender, por ejemplo,

\[
  \text{rojo}\longrightarrow\text{acción 1},
  \qquad
  \text{verde}\longrightarrow\text{acción 2}.
\]

Ya no buscamos una acción óptima global, sino una acción óptima
\textbf{dependiendo del contexto}.

Esto introduce el concepto de:

\concepto{\textbf{Política: situación $\longrightarrow$ acción}}

Los bandits contextuales están a mitad de camino entre los bandits simples y el
aprendizaje por refuerzo completo. En los bandits contextuales:

\begin{itemize}
  \item existe un estado o contexto;
  \item aprendemos qué acción tomar en cada contexto;
  \item la acción solamente afecta la recompensa inmediata.
\end{itemize}

En RL completo, además, tenemos que \textbf{la acción también afecta el estado futuro}.

\section{Conclusión}

Nos quedamos con que aprender no es solamente elegir lo que hoy parece mejor; también hay que invertir algunas acciones en conseguir información para decidir
mejor más adelante.

\end{document}